\section{Changes to the block body}\label{sec:block body changes}

The block body needs to be modified to account for Peras certificates which are included on chain to coordinate the end of a cooldown period.
We note that certificates are morally part of the Consensus layer; however, certificates are likely too large to be stored in a header.
Therefore, block bodies for Peras optionally contain a certificate next to the sequence of transactions; note that a block body then no longer contains exclusively Ledger-related content.

Concretely, we propose the following checks for a block body $B$ in round $r$ containing a certificate \cert{}:
\begin{enumerate}
\item
  The certificate \cert{} must be valid.
\item
  The round of \cert{} must be strictly greater than the round of the previous certificate on chain, and it must lie between $r-\perasCertMaxRounds$ and $r$, i.e.\
  \[ r-\perasCertMaxRounds \le \round(\cert) \le r \;. \]
  This ensures that \cert{} is neither too old nor too new, such that the state that $B$ is validated against contains the necessary information to validate $r$.
\end{enumerate}
The only change to the (extended) ledger state is the addition of the round number of the last certificate on chain, and the protocol parameters and stake distribution for the previous epoch in case \cert{} is in a previous epoch compared to $B$.%
\footnote{The ledger state already keeps this data of the previous epoch around for its reward calculation.}

\subsection{Adversarial inclusion of certificates in block body}\label{sec:block body analysis}
These rules allow the adversary to needlessly include certificates in chain even if there is no current/upcoming cooldown period.
This does not impact the purpose of on-chain certificates:
If the system does enter a cooldown (via honest nodes stopping to vote), the adversary will run out of certificates to include on chain due to the monotonicity property of round numbers. Moreover, this monotonicity property also ensures that the adversary can (on average) only include one certificate per Peras round.

We also note that implementing Peras entails diffusing \perasN{} votes per round; whose cumulative size  is significantly larger than the max size of a certificate (and the bandwidth induced by this is currently not planned to be compensated/incentivized), relativizing the actual impact of adversarial inclusion of certificates in block body.

Nodes need to validate certificates included on chain, even if they are ones that have been uselessly included by an adversary; however as it does not affect the ledger state, this validation can in principle be performed in parallel to the ledger rules, partially mitigating the impact.

Finally, the adversarially included certificates are overt and hence easy to detect for a vigilant observer of the chain, enabling \emph{social} repercussions as a partial mitigation.

Honest nodes on the other hand have no incentive to include such useless certificates, in particular as they usually could instead include more transactions (which pay fees) in their blocks.

% If this cost is considered to be unacceptable, a drastic measure would be to require pools to pay a comparable fee when they include a certificate in a block (which is a rare thing for honest nodes).
% However, such a fee for pools is unprecedented, and raises various unresolved questions regarding incentives and what to do when honest pools run out of funds to pay this new fee.

% Another option would be to modify the Peras protocol to reliably determine whether an honest node could have included a certificate into a particular block.
% However, we are not aware of an easy way to do this.


%%% Local Variables:
%%% mode: latex
%%% TeX-engine: xetex
%%% TeX-master: "../peras-design"
%%% End:
